\section{Målsetninger for markedsstrategien}
\label{sec:malsetninger}

Målene formuleres etter SMART-prinsippet, det vil si at de skal være
spesifikke, målbare, oppnåelige, relevante og tidsavgrensede
\parencite{kotlerkeller2016}. To forhold skilles ut som betingelser, ikke
markedsmål, fordi de er nødvendige forutsetninger for at markedsmålene
skal kunne nås.

\subsection{Betingelser}

\textbf{BVLOS-godkjenning innen 18 måneder.}
Luftfartstilsynets tillatelse er en juridisk forutsetning for kommersiell
drift \parencite{luftfartstilsynet2023}. Wing opererer allerede i Finland
\parencite{wing2024}, og uten tidlig regulatorisk forankring faller
førsteetableringsfortrinnet bort.

\textbf{Minst to norske innholdspartnere ved lansering.}
Verdiforslaget står og faller på tilgang til attraktive partnere. I
oppstartsfasen er Zipline mer avhengig av partnerne enn omvendt (jf.\
kraft 2 i Porters femkraftsanalyse), og partnervalget setter dermed
rammer for hvilke målgrupper planen kan adressere.

\subsection{Markedsmål}

\textbf{Oppnå en adopsjonsrate på 4 til 5 prosent i målmarkedet innen
12 måneder etter første lansering, tilsvarende rundt 500 aktive kunder
i et lokalmarked på 20\,000 innbyggere.}
Målet er spesifikt og målbart ved antall kunder med minst én gjennomført
bestilling i løpet av siste 30 dager, satt i forhold til folketallet i
målmarkedet. Oppnåeligheten støttes av Rogers' innovasjonsadopsjonsteori,
som deler en målgruppe inn etter hvor raskt de tar i bruk en ny tjeneste
\parencite{rogers2003}. Innovatørene er de første 2,5 prosent som prøver
tjenesten, og de tidlige etterfølgerne utgjør de neste 13,5 prosent. En
adopsjonsrate på 4 til 5 prosent i løpet av første driftsår tilsvarer å
fange hele innovatørgruppen og en mindre andel av de tidlige
etterfølgerne, et realistisk nivå i oppstartsåret. Den absolutte
referansen på rundt 500 aktive kunder følger av folketallet i
målmarkedet etter justering for typisk bestillingsfrekvens i nye
matleveringstjenester. Tidsrammen på 12 måneder
er valgt fordi den dekker en full sesongsyklus i hjemleveringsmarkedet
og er en konvensjonell evalueringsperiode for nye tjenester i
oppstartsåret. En kommersiell
tilstedeværelse av denne størrelsen er strategisk relevant som
forutsetning for at videre skalering skal være meningsfull.
Lavere oppnåelse tyder på at verdiforslaget eller distribusjonsmodellen
ikke har trengt gjennom i målmarkedet, og strategien må revideres.

\textbf{Etablere posisjonen «raskere og rimeligere» som foretrukket
alternativ i målmarkedet innen 12 måneder etter lansering.}
Målet måles ved en markedsundersøkelse 12 måneder etter lansering. 
Minst 15 prosent av forbrukerne i målmarkedet skal kjenne igjen 
Zipline når merket presenteres sammen med Foodora og Wolt, og 
assosiere det med raskere og rimeligere levering. Nivået ligger
innenfor det bransjebenchmarks rapporterer som indikator på
utviklende markedstilstedeværelse for nye merker i målmarked
\parencite{helmsworkshop_brandawareness}. I tillegg skal minst 50
prosent av forbrukere som har prøvd både Zipline og minst én konkurrent
oppgi Zipline som førstevalg ved tidskritiske bestillinger. Denne
terskelen er en intern KPI som måler at differensieringen oppleves som
overlegen i faktisk bruk, og er ikke en publisert bransjenorm.
Posisjonen er strategisk relevant fordi pris og hastighet er
sterke differensiatorer for Zipline (jf.\ SWOT), og uten dokumentert
posisjonering svikter forutsetningen for resten av planen. Lavere
resultat på enten posisjonsassosiasjon eller førstevalg tyder på at
verdiforslaget ikke er sterkt nok til å overstyre forbrukernes vaner, og
kommunikasjonen må revideres.

\textbf{Etablere kommersiell drift i minst to nye markeder utover første
lokalmarked innen 24 måneder etter første lansering.}
Målet er spesifikt og målbart ved at hvert nytt marked skal oppnå en
adopsjonsrate på minst 4 prosent av befolkningen i målområdet, målt som
månedlig aktive kunder, og dokumentert kontinuerlig drift i minimum 6
måneder ved utgangen av perioden. Adopsjonsraten er valgt fremfor et
absolutt kundetall slik at terskelen blir sammenlignbar på tvers av
markeder med ulik størrelse. Nivået følger samme adopsjonslogikk som i
første markedsmål \parencite{rogers2003}. Oppnåeligheten støttes av at
kjennskap opparbeidet i første lokalmarked kan redusere
adopsjonsterskelen i nærliggende områder, og av at minimumskravet på 6
måneders kontinuerlig drift gir tilstrekkelig tid til å verifisere at
adopsjonsraten er stabil. Skalerbarheten er strategisk
relevant fordi ett enkeltmarked ikke gir tilstrekkelig grunnlag for
nasjonal vekst eller dokumentasjon på at modellen tåler ulike
kundeprofiler. Lavere oppnåelse tyder på at ekspansjonsmodellen er for
kapitalintensiv eller at konkurransepresset er undervurdert, og videre
skalering må vurderes på nytt.
